\section*{第 \ref{ch:b3:lithosphere-thermal-models} 章——岩石圈热模型}

\begin{solution}{exo:b3:lithosphere-thermal-models:1}
$L^2/\kappa = \qty{e12}{s}$、\qty{e14}{s} 和 \qty{e16}{s}：0.032、3.2 和
\qty{317}{Ma}。一千米厚的岩石约在数万年内就忘掉了它的温度，
一百千米则需要数亿年。
\end{solution}

\begin{solution}{exo:b3:lithosphere-thermal-models:2}
$u = \qty{e-3}{m}/\qty{3.156e7}{s} = \qty{3.17e-11}{m.s^{-1}}$，因此
$\mathrm{Pe} = 3.17\times10^{-11}\times\num{3e4}/10^{-6} = 0.95$。对流与
传导同等重要：上升的岩石把热量向上带去，
使近地表的地温线变陡，若仅计传导，
则会低估梯度。
\end{solution}

\begin{solution}{exo:b3:lithosphere-thermal-models:3}
$\sqrt{\kappa t} = \sqrt{\num{0.8e-6}\times5\times\num{3.156e13}} =
\qty{11236}{m}$，$\eta = 10000/(2\times11236) = 0.445$，且 $T = 1450\times0.471
= \qty{683}{\celsius}$。
\end{solution}

\begin{solution}{exo:b3:lithosphere-thermal-models:4}
$490/\sqrt1 = 490$，$490/5 = 98$ 和 $490/10 = \qty{49}{mW.m^{-2}}$。探针
最可能低估第一个值：年轻的地壳裂隙发育，与
循环海水相通，部分热量通过热泉被带走。
\end{solution}

\begin{solution}{exo:b3:lithosphere-thermal-models:5}
$z = 2\times1.163\sqrt{\kappa t}$：取 $\sqrt{\kappa t} = 15.89$、31.78 和
\qty{50.25}{km}，$z = 37.0$、73.9 和 \qty{116.9}{km}。冷却层的厚度
从 10 到 \qty{40}{Ma} 翻了一番；到 \qty{100}{Ma} 时，已超过一块
\qty{95}{km} 厚的板块。
\end{solution}

\begin{solution}{exo:b3:lithosphere-thermal-models:6}
$C$ 与 $\alpha$ 成正比：$\alpha = \num{3.1e-5}\times350/370.5 =
\qty{2.93e-5}{K^{-1}}$。$C$ 也与 $T_m - T_0$ 成正比（地幔温度为
$1450\times350/370.5 = \qty{1370}{\celsius}$），并与 $\sqrt\kappa$ 成正比（
扩散率 $0.8\times(350/370.5)^2 = \qty{0.71e-6}{m^2.s^{-1}}$）：仅凭深度
无法将它们区分，这正是还要使用热流的原因。
\end{solution}

\begin{solution}{exo:b3:lithosphere-thermal-models:7}
$3.1\times1450/\num{95e3} = \qty{47.3}{mW.m^{-2}}$ 和 $3.1\times1350/\num{125e3}
= \qty{33.5}{mW.m^{-2}}$。半空间分别在 $(490/47.3)^2 =
\qty{107}{Ma}$ 和 $(490/33.5)^2 = \qty{214}{Ma}$ 时达到这两个值：较薄、较热的板块
从约 \qty{100}{Ma} 起保持老洋底的热流 \qty{47}{mW.m^{-2}}，
而较厚、较冷的板块则会使热流降至
\qty{40}{mW.m^{-2}} 以下，该值为半空间在 \qty{150}{Ma} 时的值。
\end{solution}

\begin{solution}{exo:b3:lithosphere-thermal-models:8}
$A_0h_r = 0.4\times60 = \qty{24}{mW.m^{-2}}$，因此 $A_0 =
\num{24e-3}/\num{e4} = \qty{2.4}{\micro W.m^{-3}}$，且 $q_m = 60 - 24 =
\qty{36}{mW.m^{-2}}$。在 \qty{35}{km} 深处，$T = \num{36e-3}\times\num{35e3}/2.5 +
\num{2.4e-6}\times\num{e8}/2.5\times(1 - e^{-3.5}) = 504 + 96\times0.970 =
\qty{597}{\celsius}$。
\end{solution}

\begin{solution}{exo:b3:lithosphere-thermal-models:9}
$h_r = \Delta q_0/\Delta A_0 = \num{18e-3}/\num{1.8e-6} = \qty{10}{km}$，且
$q_m = q_0 - h_rA_0 = 66 - 30 = \qty{36}{mW.m^{-2}}$，因为 $h_rA_0 =
\num{e4}\times\qty{3.0e-6}{W.m^{-3}} = \qty{30}{mW.m^{-2}}$。
\end{solution}

\begin{solution}{exo:b3:lithosphere-thermal-models:10}
(a) $Q(t) = \int_0^tk(T_m - T_0)/\sqrt{\pi\kappa t'}\,\mathrm dt' =
2k(T_m - T_0)\sqrt{t/(\pi\kappa)}$。(b) 散失的热量即是从柱体中移走的热含量，
$Q = \rho_mc\int_0^\infty(T_m - T)\,\mathrm dz$，而由
均衡条件得 $(\rho_m - \rho_w)(d - d_r) = \rho_m\alpha\int_0^\infty(T_m -
T)\,\mathrm dz = \alpha Q/c$。(c) $c = 3.1/(3300\times\num{0.8e-6}) =
\qty{1174}{J.kg^{-1}.K^{-1}}$；在 \qty{40}{Ma} 时，$t = \qty{1.262e15}{s}$，且
$Q = 2\times3.1\times1450\times\sqrt{\num{1.262e15}/(\pi\times\num{0.8e-6})} =
\qty{2.015e14}{J.m^{-2}}$；$d - d_r = \num{3.1e-5}\times\num{2.015e14}/
(1174\times2270) = \qty{2343}{m}$，即 $370.5\sqrt{40}$。
\end{solution}

\begin{solution}{exo:b3:lithosphere-thermal-models:11}
第 $n = 3$ 项为 $3104\times(8/\pi^2)\times\tfrac19e^{-9t/36.2} =
280\,e^{-9t/36.2}$ 米：在 \qty{1}{Ma} 时为 \qty{218}{m}，在
\qty{10}{Ma} 时为 \qty{23}{m}，\qty{20}{Ma} 时为 \qty{1.9}{m}。当其降至 \qty{10}{m} 时，
$e^{-9t/36.2} = 10/280$ 时，$t = (36.2/9)\ln28 = \qty{13.4}{Ma}$。
\end{solution}

\begin{solution}{exo:b3:lithosphere-thermal-models:12}
(a) 当 $z \gg h_r$ 时，$e^{-z/h_r} \to 0$。(b) $q_m = \qty{24}{mW.m^{-2}}$，
$A_0h_r^2/k = \num{16e-3}\times\num{e4}/2.5 = \qty{64}{K}$，且 $z = (1300 -
64)\times2.5/\num{24e-3} = \qty{128.8}{km}$。(c) $q_m = 2.5\times1236/\num{250e3}
= \qty{12.4}{mW.m^{-2}}$：\omterm{def:b3:lithosphere-thermal-models:reduced}{折合热流}须约减半，如同克拉通之下的情形，
其厚而冷的根部从下方导入的热量很少。
\end{solution}

\begin{solution}{pb:b3:lithosphere-thermal-models:1}
\textbf{1.} $\sqrt{\num{0.8e-6}\times\num{3.156e14}} = \qty{15889}{m}$ 和
$\sqrt{\num{0.8e-6}\times\num{1.2624e15}} = \qty{31778}{m}$。

\textbf{2.} $\eta = 20000/31778 = 0.6294$ 和 $20000/63556 = 0.3147$：
$T = 1450\times0.6266 = \qty{909}{\celsius}$（\qty{10}{Ma} 时），且
$1450\times0.3437 = \qty{498}{\celsius}$（\qty{40}{Ma} 时）。

\textbf{3.} $z = 2\times0.7174\sqrt{\kappa t}$：\qty{22.8}{km} 和
\qty{45.6}{km}。深度达到 \qty{60}{km} 时，$\sqrt{\kappa t} = 60/1.4348 =
\qty{41.82}{km}$，此时 $t = (\num{41.82e3})^2/\num{0.8e-6} = \qty{2.186e15}{s}$，
即 \qty{69.3}{Ma}。

\textbf{4.} $q = 3.1\times1450/(\sqrt\pi\times15889) = \qty{159.6}{mW.m^{-2}}$（\qty{10}{Ma} 时）；
\qty{79.8}{mW.m^{-2}} 为 \qty{40}{Ma} 时的值；$C_q = 159.6\sqrt{10}
= \qty{505}{mW.m^{-2}.Ma^{1/2}}$。

\textbf{5.} $u = \qty{7.92e-10}{m.s^{-1}}$，$\mathrm{Pe} = 7.92\times10^{-10}
\times\num{95e3}/\num{0.8e-6} = 94$。热量随板块水平输运，
速度比水平扩散快约一百倍，因此每个柱体
仅靠竖直传导冷却。

\textbf{6.} $(\num{95e3})^2/\num{0.8e-6} = \qty{1.13e16}{s}$，即 \qty{357}{Ma}。
当 \qty{150}{Ma} 时，$\sqrt{\kappa t} = \qty{61.5}{km}$，冷却层
厚度为其若干倍，已远深于 \qty{95}{km}：对如此古老的洋底采用
半空间模型，意味着冷的影响已达到地幔对流的深度。

\textbf{7.} $(\rho_m - \rho_w)(d - d_r) = \rho_m\alpha\int(T_m - T)\,\mathrm dz$，
在补偿深度以上，洋脊柱体与老洋底柱体的
质量相等。$T_m - T = T_m\operatorname{erfc}\eta$ 降至
$T_m$ 的 \qty{0.5}{\percent} 以下，对 $\eta > 2$ 成立，因此深部对结果没有贡献。

\textbf{8.} $C = 2\times3300\times\num{3.1e-5}\times1450/2270\times
\sqrt{\num{0.8e-6}\times\num{3.156e13}/\pi} = 0.13069\times2834.8 =
\qty{370.5}{m.Ma^{-1/2}}$；$d = 2500 + 370.5\sqrt{10} = \qty{3672}{m}$，和
$2500 + 370.5\sqrt{40} = \qty{4843}{m}$。

\textbf{9.} 模型给出 $2500 + 370.5\sqrt{42.5} = \qty{4915}{m}$，
比中位数深 \qty{440}{m}。中位数包含沉积物的负荷、
海山和火山高原、热点隆起，以及
断裂带和大陆边缘两侧的对流环流，这些都使洋底
变浅。

\textbf{10.} $C$ 与 $T_m$ 成正比：$370.5\times1350/1450 =
\qty{344.9}{m.Ma^{-1/2}}$，与拟合值 350 接近。

\textbf{11.} 海水在破裂的年轻地壳中循环，
在探针之间通过热泉带走部分热量。

\textbf{12.} $\tau = (\num{95e3})^2/(\pi^2\times\num{0.8e-6}) =
\qty{1.143e15}{s} = \qty{36.2}{Ma}$；振幅为 $3300\times
\num{3.1e-5}\times1450\times\num{95e3}/4540 = \qty{3104}{m}$。

\textbf{13.} $e^{-40/36.2} = 0.33124$，$e^{-360/36.2}/9 = \num{5.3e-6}$，而
$n = 5$ 项更小；括号内为 $1 - 0.81057\times0.33125 =
0.73150$，且 $d = 2500 + 3104\times0.73150 = \qty{4771}{m}$。

\textbf{14.} 由 $n = 3$ 项可得 $3104\times0.81057\times\num{5.3e-6} =
\qty{0.013}{m}$：在 \qty{40}{Ma} 时，$e^{-9t/\tau}$ 约为
$10^{-4}$ 倍的 $e^{-t/\tau}$，只剩第一模态。

\textbf{15.} 在 \qty{40}{Ma} 时，$4843 - 4771 = \qty{73}{m}$；在
\qty{50}{Ma} 时，$2500 + 370.5\sqrt{50} = 5120$，而 $2500 + 3104\times(1 -
0.81057\times0.25127) = 4972$，两者相差 \qty{148}{m}。线性内插得 $40 + 10\times
(100 - 73)/(148 - 73) = \qty{43.6}{Ma}$，约为 \qty{44}{Ma}（完整级数
给出 44.0）。

\textbf{16.} $q = (3.1\times1450/\num{95e3})[1 + 2(0.33124 + 0.01204 +
0.00005)] = 47.32\times1.68666 = \qty{79.8}{mW.m^{-2}}$，与半空间模型结果相同：
半空间模型在 \qty{40}{Ma} 时，底部的影响尚未
传到地表。

\textbf{17.} $2500 + 3104 = \qty{5604}{m}$ 和 $kT_m/L = \qty{47.3}{mW.m^{-2}}$。

\textbf{18.} $d(120) = 2500 + aL\bigl[1 - (8/\pi^2)e^{-120/\tau}\bigr]$，其中
$a = \rho_m\alpha T_m/\bigl(2(\rho_m - \rho_w)\bigr) = \qty{32.673}{m/km}$，且
$\tau = L^2/(\pi^2\kappa)$。

\textbf{19.} $L = 3200/\bigl(32.673\,[1 - 0.81057e^{-120/\tau}]\bigr)$。由
$L = \qty{100}{km}$ 出发：$\tau = \qty{40.13}{Ma}$，括号项为 0.95925，$L =
\qty{102.10}{km}$；然后 $\tau = 41.83$，$L = 102.67$；$\tau = 42.30$，$L =
102.82$；$\tau = 42.43$，$L = 102.87$；最终收敛于 $L = \qty{102.9}{km}$。

\textbf{20.} $\tau = \qty{42.48}{Ma}$，$aL = \qty{3362}{m}$：$d(40) = 2500 +
3362\times(1 - 0.81057e^{-40/42.48}) = \qty{4799}{m}$，与半空间的 \qty{4843}{m}
以及实测的 \qty{4800}{m} 相比：该板块与两个年龄的数据都相符。

\textbf{21.} $\tau = \qty{42.5}{Ma}$；渐近深度为 $2500 + 3362 =
\qty{5862}{m}$。

\textbf{22.} $3362\times0.81057\,e^{-t/42.48} = 100$ 给出 $t = 42.48\ln27.25 =
\qty{140}{Ma}$。

\textbf{23.} 与虚构大洋老洋底相吻合的板块，其厚度为
\textbf{\qty[text-series-to-math=true]{102.9}{km}}。
\end{solution}
